% ═══════════════════════════════════════════════════════════════════════════
% QR-CODE: Simulation Oxidation (Chemie Klasse 8)
% ═══════════════════════════════════════════════════════════════════════════

\documentclass[a4paper,11pt]{article}

\usepackage[utf8]{inputenc}
\usepackage[T1]{fontenc}
\usepackage{helvet}
\renewcommand{\familydefault}{\sfdefault}

\usepackage[margin=20mm]{geometry}
\usepackage{tikz}
\usepackage{xcolor}
\usepackage{qrcode}

% ═══════════════════════════════════════════════════════════════════════════
% FARBEN
% ═══════════════════════════════════════════════════════════════════════════

\definecolor{chemie}{HTML}{2a7a4b}
\colorlet{fachfarbe}{chemie}

% ═══════════════════════════════════════════════════════════════════════════
% KONFIGURATION
% ═══════════════════════════════════════════════════════════════════════════

\newcommand{\simtitel}{Oxidation von Metallen}
\newcommand{\simurl}{https://devbox42.github.io/sciencesim/chemie/kl08/01-metalloxide/SIM-01-oxidation.html}
\newcommand{\simurlkurz}{devbox42.github.io/.../SIM-01-oxidation.html}

% ═══════════════════════════════════════════════════════════════════════════
% DOKUMENT
% ═══════════════════════════════════════════════════════════════════════════

\pagestyle{empty}

\begin{document}

\begin{center}

\vspace*{10mm}
{\fontsize{28}{34}\selectfont\bfseries\color{fachfarbe}Simulation: \simtitel}

\vspace{15mm}

\fcolorbox{fachfarbe}{white}{\qrcode[height=100mm]{\simurl}}

\vspace{12mm}

{\large\color{gray}\simurlkurz}

\vspace{8mm}

{\Large\color{gray}Scanne den QR-Code mit deinem Gerät}

\vspace{10mm}

\fcolorbox{fachfarbe}{fachfarbe!10}{%
    \parbox{0.7\textwidth}{%
        \centering
        \vspace{3mm}
        Wähle ein Metall und beobachte die Reaktion mit Sauerstoff.\\[2mm]
        Notiere Flammenfarben und Oxidfarben auf dem Arbeitsblatt.
        \vspace{3mm}
    }%
}

\end{center}

\end{document}
